\documentclass[12pt]{article}
\usepackage[utf8]{inputenc}
\usepackage[T1]{fontenc}
\usepackage{hyperref}
\title{Teduray Language Resources for K-12 Curriculum Development}
\author{Compiled by Claude Sonnet 5 (claude-sonnet-5) for Andr\'as Kornai}
\date{}
\begin{document}
\maketitle
\section{Introduction}
Teduray (Tiruray, ISO 639-3 tiy) is an Austronesian language of Maguindanao and Sultan Kudarat provinces in the Philippines. The searches found a Bible translation and general descriptions. They found no dictionary, grammar or school materials.

\section{Primary Educational Resources}
No primers, textbooks or mother-tongue education programme for Teduray were found.

\section{Linguistic and Reference Materials}
\begin{itemize}
\item Wikipedia on the Teduray language, with its phonology and structure: \url{https://en.wikipedia.org/wiki/Teduray_language}
\item Omniglot on Tiruray, which describes the Latin-alphabet writing: \url{https://omniglot.com/writing/tiruray.htm}
\end{itemize}

\section{Religious and Community Resources}
ScriptureEarth lists a Teduray Bible, with a New Testament PDF and ePub, and a Jesus Film: \url{https://www.scriptureearth.org/00eng.php?iso=tiy}

\section{Digital and Audio Resources}
The ScriptureEarth page mentions a 2002 Wycliffe audio recording on Bible.is. No other audio resources were found.

\section{Curriculum Development Context}
The Bible text is the only substantial written material located. A curriculum would need new materials, built with the community and drawing on any SIL or Philippine mother-tongue education work, which was not found here.
\end{document}